% Teil II: der durchgerechnete Weg.

\subsection*{Gegeben und gesucht}
Gegeben: $R_E=\qty{6378}{km}=\qty{6.378e6}{m}$, $M_E=\qty{5.974e24}{kg}$,
$G=\qty{6.674e-11}{m^3/(kg\,s^2)}$.
Gesucht: anfängliche Geschwindigkeit $v_0$ zum Verlassen des Gravitationsfelds.

\subsection*{Ansatz}
Energieerhaltung zwischen Oberfläche (Zustand 1) und Unendlich (Zustand 2).
Im Unendlichen soll $v_2=0$ sein (gerade noch entkommen). Damit die Masse
$m$ der Kugel aus der Gleichung.

\subsection*{Rechnung}
\begin{equation}
  \frac12 m v_0^2 - G\frac{M_E m}{R_E}
  = \frac12 m\,0^2 - 0
  = 0\;.
\end{equation}
Nach $v_0$ auflösen ($m$ kürzt sich):
\begin{equation}
  v_0 = \sqrt{\frac{2\,G M_E}{R_E}}
  = \sqrt{\frac{2\cdot 6.674e-11\cdot 5.974e24}{6.378e6}}\ \si{m/s}
  \approx \qty{1.12e4}{m/s}
  \approx \qty{11.2}{km/s}.
\end{equation}

\begin{ergebnis}
\begin{equation*}
  v_0 = \sqrt{\frac{2GM_E}{R_E}} \approx \qty{11.2}{km/s}.
\end{equation*}
Die Kugel muss (reibungslos, vertikal) mindestens \qty{11.2}{km/s} schnell sein,
um das Gravitationspotential der Erde zu verlassen.
\end{ergebnis}

\begin{kontrolle}
\textbf{Einheiten:} $G M_E/R_E$:
$\frac{m^3}{kg\,s^2}\cdot\frac{kg}{m}=m^2/s^2$, damit $\sqrt{\cdots}=m/s$. \checkmark

\textbf{Grössenordnung:} Die Fluchtgeschwindigkeit von der Erde ($\approx
\qty{11.2}{km/s}$) liegt um $\sqrt2$ über der ersten kosmischen
Geschwindigkeit für eine erdnahe Kreisbahn ($\approx\qty{7.9}{km/s}$);
das ist physikalisch bekannt (Kreisbahn braucht $\sqrt{GM_E/R_E}$).

\textbf{Grenzfall:} Bei $R_E\to\infty$ (keine Masse) würde $v_0\to 0$.
\end{kontrolle}

\begin{merksatz}
Entkommgeschwindigkeit = $\sqrt{2GM/R}$: die Energie wird von der
potentiellen auf die kinetische Form übertragen; Masse des fliegenden
Objekts kürzt sich heraus. Gilt analog für jede Masse im Gravitationsfeld
einer kugelsymmetrischen Verteilung.
\end{merksatz}
